\section{Introduction}\label{sec:introduction}

Let $K$ be a smooth knot in the oriented three-sphere, and let
$E_K=S^3\setminus\operatorname{int}(\nu K)$ be its exterior.
The meridian $\mu$ and preferred longitude $\lambda$ identify the unoriented
slopes on $\partial E_K$ with $\Q\cup\{\infty\}$: the primitive class
$p\mu+q\lambda$, considered up to simultaneous sign, represents $p/q$.
Write $S^3_{p/q}(K)$ for the corresponding Dehn filling, with the orientation
extending that of $E_K$.
Two distinct slopes giving orientation-preservingly homeomorphic fillings
are called \emph{purely cosmetic}.

\begin{theorem}\label{thm:main}
Let $K\subset S^3$ be a smooth nontrivial knot. For any
$r,s\in\Q\cup\{\infty\}$, an orientation-preserving homeomorphism
\[
 S^3_r(K)\cong S^3_s(K)
\]
implies $r=s$.
\end{theorem}

Theorem~\ref{thm:main} resolves the purely cosmetic surgery conjecture for
knots in the three-sphere. The homeomorphism in its statement is a map of
the entire filled manifolds; it need not preserve a surgery core or the
knot exterior. The orientation hypothesis is essential to the question:
orientation-reversing, or \emph{chirally cosmetic}, coincidences are a
different phenomenon.

\subsection{The remaining case}

Dehn surgery converts the peripheral information of a knot into a closed
three-manifold. The cosmetic surgery problem asks how much of the filling
slope can be recovered from that oriented manifold. For knot exteriors,
this is an instance of the cosmetic filling problem formulated by
Gordon~\cite[Conjecture~6.1]{Gordon1991} and recorded as Bleiler's
Problem~1.81(A) in Kirby's list~\cite{Kirby1997}.
Bleiler, Hodgson, and Weeks developed the terminology and distinguished
fillings related by a symmetry of the exterior from those that are
not~\cite{BleilerHodgsonWeeks1999}.  Here the homeomorphism between
fillings is arbitrary.  Gordon and Luecke's theorem that nontrivial
surgery on a nontrivial knot cannot produce $S^3$
\cite[Theorem~2]{GordonLuecke1989} settles the meridional case.

Several theories have supplied obstructions to a cosmetic pair of finite slopes.
Boyer and Lines obtained an Alexander-polynomial obstruction from
surgery formulae for classical three-manifold invariants~\cite{BL90}.
The rational surgery formula of Ozsv\'ath and Szab\'o~\cite{OS11},
Wang's genus-one theorem~\cite{Wang06}, and Wu's exclusion of
same-sign purely cosmetic pairs~\cite{Wu11} established a Heegaard
Floer approach.  Ni and Wu strengthened the restriction to opposite
slopes~\cite[Theorem~1.2]{NiWu2015}.  Using graded Floer information
through immersed curves, Hanselman restricted the possibilities to
$\pm2$ or $\pm1/q$ for a positive integer $q$, and showed that the
$\pm2$ case requires genus two~\cite[Theorem~2]{Hanselman2023}.

Other obstructions and special cases clarify the range of the problem.
Ichihara and Wu used the Jones polynomial~\cite{IchiharaWu2019},
and Detcherry developed a quantum obstruction~\cite{Detcherry2026}.
Tao settled nontrivial cable and composite knots
\cite{TaoCables2019,TaoConnected2022}, and Stipsicz and Szab\'o
settled nontrivial pretzel knots~\cite{StipsiczSzabo2021}.
Kotelskiy, Lidman, Moore, Watson, and Zibrowius proved an equivariant
version for strongly invertible knots, in which the diffeomorphism
commutes with the extended involutions
\cite{KotelskiyLidmanMooreWatsonZibrowius2024}.
Futer, Purcell, and Schleimer combined knot invariants, hyperbolic
estimates, and verified computations to establish the conjecture
through $19$ crossings and for the one-cusped SnapPy census
\cite{FuterPurcellSchleimer2025}.  Their procedure applies beyond
these finite data sets but is not guaranteed to distinguish every pair.
Ren proved that the conjecture is equivalent to its restriction to
hyperbolic surgeries on hyperbolic knots
\cite[Theorem~3, arXiv version~2]{Ren2025}.

The filtered instanton argument of Daemi, Lidman, and Miller Eismeier
excludes the pairs $\pm1/q$.  Combined with the earlier restrictions,
their Corollary~1.4 leaves only opposite integer slopes of magnitude
two on a genus-two knot; it also imposes additional Alexander and Jones
polynomial conditions~\cite{DaemiLidmanMillerEismeier2025}.
We use only the following consequence.

\begin{proposition}[Cosmetic surgery reduction]\label{prop:introduction:reduction}
If a smooth nontrivial knot $K\subset S^3$ has distinct purely cosmetic
slopes, then $g(K)=2$ and the two slopes are $-2$ and $2$.
\end{proposition}

For completeness, the zero slope causes no additional case:
$H_1(S^3_{p/q}(K);\Z)$ is cyclic of order $|p|$ for $p\ne0$ and is
infinite cyclic for $p=0$. Primitivity makes $0/1$ the unique slope with
zero numerator. The meridional slope is covered by the sphere-filling
theorem just mentioned.

We henceforth suppose that $g(K)=2$ and that there is an oriented
diffeomorphism
\begin{equation}\label{eq:introduction:cosmetic}
 \phi:Y_{-2}\longrightarrow Y_2,
 \qquad Y_j=S^3_j(K),\quad Y_\infty=S^3.
\end{equation}
The three-dimensional smoothing theorem permits the original
homeomorphism to be replaced by a homotopic diffeomorphism
\cite[Theorem~2.3]{Massuyeau2011}.  Homotopy preserves its degree,
so this replacement preserves orientation.  We will obtain a
contradiction without imposing any peripheral condition on $\phi$.

\subsection{The argument}

There are two counts. The first is an integer instanton count
$\Omega\ne0$ on a closed parameterized cobordism. The second is an
oriented one-dimensional spinor count whose boundary would force
$2^\eta\Omega=0$. The work is to construct the two counts with compatible
metrics, incidence conditions, determinant lifts, and orientations.

The comparison follows the nonabelian-monopole cobordism strategy of
Pidstrigach--Tyurin and Okonek--Teleman
\cite{PidstrigachTyurin1995,OkonekTeleman1996}, with the analytic and
instanton-link development of Feehan--Leness
\cite{FeehanLenessPU2I,FeehanLenessPU2II,FeehanLenessLinksI}.
The doubled phase class is part of that method.  Here the simultaneous
family, its relative boundary analysis, and the regularization of
mixed projections are constructed explicitly.

The ordinary count uses instanton Floer and Donaldson theory
\cite{Floer1988,Donaldson2002,DonaldsonKronheimer1990}, with
determinant-one gauge transformations.
At the rational ends with first homology of order at most two, all flat
reducers are central. At zero surgery we instead use the determinant
class evaluating oddly on the capped Seifert surface. Ordinary surgery
families give units modulo two. The negative direction requires an
interval of metrics, a degree-two sphere test, and a signed sum over two
determinant lifts. The two minimal lens-wall contributions cancel with
these signs. Symplectic caps and ordinary Donaldson invariants provide a
nonzero rational cap pairing whose probes can be moved off all variable
blocks, following the cap and nonvanishing methods of
Kronheimer--Mrowka~\cite{KronheimerMrowka2004,KronheimerMrowka1995}.
Repetition on the free integral Floer lattice modulo a power of
two then produces the nonzero integer $\Omega$.

The same repetition makes the complex coupled Dirac index $n_D$
positive while preserving nonvanishing; the index calculation and
integral congruences are in Section~\ref{sec:stack}.
The phase circle acts by multiplying the spinor.  Starting from the
zero-dimensional ordinary count, adding spinors, imposing
$\eta=n_D-1$ complex phase cuts, and taking the effective phase
quotient gives expected dimension
\[
 2n_D-2\eta-1=1.
\]
Near an ordinary instanton the phase link is $\CP^{n_D-1}$, and each
phase cut has class twice the hyperplane class, as in the link
calculation of Feehan--Leness~\cite[Lemma~3.28 and
Proposition~3.29, arXiv version~3]{FeehanLenessPU2II}.  These links contribute
$2^\eta\Omega$ with the ordinary instanton signs.

Reducible fields and broken or bubbled limits could supply additional
ends.  To control them, we construct toric metric sectors using the
Hessian description of Guillemin and the scalar-curvature formula of
Abreu~\cite{Guillemin1994,Abreu1998}.  Their periods constrain every
parallel line reduction, that is, a splitting of the bundle into line
bundles preserved by the connection. Long tube estimates and
isolated-cap estimates make those constraints simultaneous, including
at genuine corners of the parameter family. The resulting quadratic
and index inequalities exclude fixed tuples and every mixed broken or
bubbled boundary except a face at which one minimal trace cap separates.

At each surviving lens face, precisely one minimal trace cap separates
from a single free main component. These limits have no particles, and
regular cylindrical gluing applies. The two determinant sectors at each face cancel using
the same local orientation comparison that defined the negative
instanton unit. The oriented boundary formula now contradicts
$\Omega\ne0$.

The analytic construction proceeds in stages. The ordinary representatives
and their compactness properties are established independently of the
nonvanishing argument. For the spinor problem, energy bounds and ideal
compactness precede fixed-type regularization and its index exclusions;
these precede the finite projected regularization and the construction of
the remaining regular ends.  Gluing these ends precedes the final
interior regularization of the one-dimensional free stratum.
The generic choices are made around the
ordinary data, within a neighborhood preserving $\Omega$. The final
boundary formula is applied only after these stages are complete.

\begin{figure}[htbp]
\centering
\begin{tikzpicture}[
 box/.style={draw,rounded corners=2pt,align=center,text width=5.4cm,
             minimum height=1.05cm,inner sep=5pt,font=\small},
 arrow/.style={-{Stealth[length=2mm]},thick},
 node distance=9mm and 8mm]
 \node[box] (ordinary) {Ordinary units and cap pairing\\
   Sections~\ref{sec:ordinary}--\ref{sec:caps}};
 \node[box,right=of ordinary] (metrics) {Metric families and line estimates\\
   Sections~\ref{sec:geometry}--\ref{sec:estimates}};
 \node[box,below=of ordinary] (count) {Integral repetition\\
   $\Omega\ne0$, $n_D>0$; Section~\ref{sec:stack}};
 \node[box,below=of metrics] (bounds) {Boundary index inequalities\\
   Section~\ref{sec:indices}\\
   \footnotesize Hypotheses supplied in
   Sections~\ref{sec:analysis}--\ref{sec:gluing}};
 \node[box,below=16mm of $(count)!0.5!(bounds)$,text width=8.8cm] (analytic)
   {Compatible perturbations, \mbox{compactness}, and regular end gluing\\
    Sections~\ref{sec:analysis}--\ref{sec:gluing}};
 \node[below=8mm of analytic,font=\small] (boundary)
   {Instanton links: $2^\eta\Omega$ \qquad Lens pairs: $0$};
 \draw[arrow] (ordinary)--(count);
 \draw[arrow] (metrics)--(bounds);
 \draw[arrow] (count.south)--(analytic.north west);
 \draw[arrow] (bounds.south)--(analytic.north east);
 \draw[arrow] (analytic)--(boundary);
\end{tikzpicture}
\caption{The two parts of the proof meet in one compatible boundary
count. The analytic construction supplies the regularity hypotheses used
by the dimension inequalities; those inequalities in turn identify the
allowed ends.}
\label{fig:introduction:structure}
\end{figure}

\subsection{Local constructions and organization}

Three constructions carry the main technical content. First, a
parameterized negative cobordism is made into an integral chain map by a
local comparison of two determinant lifts, while its invertibility is
detected modulo two (Theorem~\ref{thm:negative:unit} and
Proposition~\ref{prop:negative:local-sign}). Second, a family of toric periods, together with a
separate flux estimate on a long $S^2\times S^1$ tube, yields line-class
inequalities that persist across compatible separations
(Theorems~\ref{thm:geometry:family} and \ref{thm:estimates:clamp}). Third,
exact-support incidence conditions and finite equivariant branch
perturbations permit negative-dimension arguments on projections of
mixed strata (Lemmas~\ref{lem:analysis:incidence} and
\ref{lem:analysis:projection}, and Proposition~\ref{prop:analysis:finite-induction}).
These projections forget zero-spinor reducer fields while retaining
their parameters; the remaining equations are independent of the
forgotten fields. The finite construction is confined to images of
candidate tuples whose retained main fields solve the equations and
cuts, with omitted lens caps and cylinder levels recorded through discrete
nonnegative charge assignments.
Their statements retain the local hypotheses explicitly,
so that these constructions can be used apart from the final knot
surgery contradiction.

The proof distinguishes two kinds of neck throughout. The surgery
seams $Y_j$ and the auxiliary isolation collars are long but finite.
The genuine infinite faces in the spinor compactification are the
specified $S^3$ cuts, denoted $J$, and the lens boundaries of
square-$-4$ sphere neighborhoods. Their limiting metrics have positive
scalar curvature. Thus the proof does not require a general regular
spinor Floer theory on the knot surgeries, or gluing at positive particle
levels.

Sections~\ref{sec:ordinary}--\ref{sec:stack} construct the ordinary
nonzero count and its topological data.
Sections~\ref{sec:geometry}--\ref{sec:estimates} construct the metrics and
prove the line estimates with their order of choices.
Section~\ref{sec:indices} proves the boundary inequalities under precise
regularity and incidence hypotheses.
Section~\ref{sec:analysis} constructs representatives, proves compactness,
and supplies the regularity needed for the boundary exclusions and the
instanton links. Section~\ref{sec:gluing} constructs the surviving lens
collars, regularizes the remaining free interior, and completes the
boundary count proving Theorem~\ref{thm:main}.
